\section{The restricted native Gram of the missing carrier}
\label{sec:tpc61-restricted-gram}

The ambient native degree from TPC--60 is the norm of a grouping operator
on a free atomic space \citep{WangTPC60}.  The physical missing vector does
not range over that free space: every atom in one source row is driven by
the same row coefficient.  This section computes the optimal grouping
constant on that actual coefficient image.  It also separates the
same-residual determinant-ladder collisions from the new collisions created
when the residual label is erased.

\subsection{Atomic, residual, and native maps}

Let \(\cW_M\) be the fixed missing-high carrier, let \(\cM_X\) be the
declared finite source-row set, and put
\[
 \cM_M:=\{m\in\cM_X:
 \text{some }w\in\cW_M\text{ satisfies }m(w)=m\}.
\]
All maps below act on the original coefficient space
\(\cK_X=\ell^2(\cM_X)\); rows outside \(\cM_M\) remain zero directions of
the missing diagonal rather than being deleted from the represented map.
On an active missing atom
\(w\), write the literal coefficient as
\begin{equation}
 b_w(\zeta)=\beta_w\zeta_{m(w)},
 \qquad
 m(w)\in\cM_M,
 \label{eq:tpc61-literal-missing-coefficient}
\end{equation}
and absorb the physical high-character sign and the residual M\"obius gauge
into
\begin{equation}
 c_w:=-\mu(s(w))\beta_w.
 \label{eq:tpc61-sign-compatible-atom}
\end{equation}
The active squarefree support gives
\(|c_w|=|\beta_w|\).  Define the atomic lift
\begin{equation}
 L_M:\cK_X\longrightarrow\ell^2(\cW_M),
 \qquad
 (L_M\zeta)_w=c_w\zeta_{m(w)}.
 \label{eq:tpc61-atomic-lift}
\end{equation}
Its Gram is diagonal:
\begin{align}
 \mathbf D_M
 &:=L_M^*L_M
   =\operatorname{diag}\bigl(W_M(m)\bigr)_{m\in\cM_X},
 \label{eq:tpc61-missing-diagonal-matrix}\\
 W_M(m)
 &:=\sum_{\substack{w\in\cW_M\\m(w)=m}}|\beta_w|^2,
 \qquad
 \cD_M(\zeta)=\langle\mathbf D_M\zeta,\zeta\rangle.
 \label{eq:tpc61-row-missing-weight}
\end{align}
Thus \(\cD_M\) is exactly the raw missing atomic diagonal, with no fiber
normalization.

Let \(i(w)\) be the native output key and \(s(w)\) the residual label.
The first grouping and the second erasure are
\begin{align}
 (T_M\zeta)_{i,s}
 &:=
 \sum_{\substack{w\in\cW_M\\i(w)=i,\ s(w)=s}}
 c_w\zeta_{m(w)},
 \label{eq:tpc61-residual-map}\\
 (\mathsf E z)_i&:=\sum_s z_{i,s}.
 \label{eq:tpc61-native-erasure}
\end{align}
All sums use the active labels of the declared carrier.  In particular,
\(\mathsf E\) does not adjoin ambient residual labels with zero physical
coefficient.

\begin{proposition}[Literal three-level factorization]
\label{prop:tpc61-three-level-factorization}
The physical missing map of TPC--60 satisfies
\begin{equation}
 \boxed{M_X=\mathsf E T_M.}
 \label{eq:tpc61-missing-factorization}
\end{equation}
Moreover \(\Ker\mathbf D_M\subseteq\Ker T_M\subseteq\Ker M_X\).
\end{proposition}

\begin{proof}
At a native coordinate \(i\), the right side of
\eqref{eq:tpc61-missing-factorization} is
\[
 \sum_s\sum_{\substack{w\in\cW_M\\i(w)=i,\ s(w)=s}}
 -\mu(s(w))\beta_w\zeta_{m(w)}.
\]
This is exactly the sign-compatible missing-high contribution in the
TPC--58/TPC--60 convention.  If
\(\zeta\in\Ker\mathbf D_M\), then every atom driven by a nonzero coordinate
of \(\zeta\) has zero fixed weight, so \(L_M\zeta=0\), and hence both
grouped maps vanish.
\end{proof}

Put
\begin{equation}
 S_M:=\Ran\mathbf D_M=(\Ker\mathbf D_M)^\perp.
 \label{eq:tpc61-active-row-space}
\end{equation}
Only the identically inactive rows are removed in \(S_M\).  In particular,
one must not quotient by \(\Ker M_X\): a nonzero vector in that kernel is a
genuine physical cancellation direction relative to the nonzero atomic
norm \(\cD_M\).

All uniform suprema below are unchanged if they are restricted to the
declared bounded coefficient cube, provided that cube contains a nonzero
scalar multiple of every direction in \(\cK_X\).  This follows from the
degree-zero homogeneity of the Rayleigh quotients; it does not grant
independent atom coefficients inside one source row.

\subsection{The optimal effective native degree}

\begin{theorem}[Exact restricted native degree]
\label{thm:tpc61-optimal-kappa}
If \(S_M\ne\{0\}\), define
\begin{equation}
 \boxed{
 \kappa_M:=
 \lambda_{\max}\!\left(
  \mathbf D_M^{\dagger/2}M_X^*M_X
  \mathbf D_M^{\dagger/2}\big|_{S_M}
 \right).}
 \label{eq:tpc61-kappa-definition}
\end{equation}
Then
\begin{align}
 \kappa_M
 &=\sup_{\cD_M(\zeta)>0}
   \frac{\|M_X\zeta\|^2}{\cD_M(\zeta)}
 \label{eq:tpc61-kappa-rayleigh}\\
 &=\inf\{\kappa\ge0:
       M_X^*M_X\preceq\kappa\mathbf D_M\}.
 \label{eq:tpc61-kappa-loewner}
\end{align}
Consequently
\begin{equation}
 \boxed{
 \|M_X\zeta\|^2\le\kappa_M\cD_M(\zeta)}
 \qquad(\zeta\in\cK_X),
 \label{eq:tpc61-optimal-native-upper}
\end{equation}
and \(\kappa_M\) is the least possible coefficient-uniform constant in
\eqref{eq:tpc61-optimal-native-upper}.  If \(S_M=\{0\}\), then
\(M_X=0\), and we set \(\kappa_M=0\).
\end{theorem}

\begin{proof}
For \(\cD_M(\zeta)>0\), put
\(x=\mathbf D_M^{1/2}\zeta\in S_M\).  Since \(M_X\) vanishes on
\(\Ker\mathbf D_M\),
\[
 M_X\zeta=M_X\mathbf D_M^{\dagger/2}x.
\]
The quotient in \eqref{eq:tpc61-kappa-rayleigh} is therefore the Rayleigh
quotient of the positive operator in
\eqref{eq:tpc61-kappa-definition}.  Rayleigh--Ritz proves
\eqref{eq:tpc61-kappa-rayleigh}.  A quadratic-form inequality for every
\(\zeta\) is equivalent to the Loewner inequality in
\eqref{eq:tpc61-kappa-loewner}, which proves optimality.
\end{proof}

Let \(C_M:\ell^2(\cW_M)\to\ell^2(\cI_{\rm out})\) be the ambient native
grouping
\begin{equation}
 (C_Ma)_i=\sum_{\substack{w\in\cW_M\\i(w)=i}}a_w,
 \qquad
 d_{\rm miss}:=\max_i|\{w\in\cW_M:i(w)=i\}|.
 \label{eq:tpc61-ambient-native-grouping}
\end{equation}
Then \(M_X=C_ML_M\) after the unit gauges have been included in
\eqref{eq:tpc61-sign-compatible-atom}.

\begin{corollary}[Compression improves the ambient degree]
\label{cor:tpc61-kappa-below-degree}
One has
\begin{equation}
 \boxed{\kappa_M\le d_{\rm miss}.}
 \label{eq:tpc61-kappa-ambient-bound}
\end{equation}
The inequality can be sharp, or the two sides can differ by an arbitrarily
large factor.
\end{corollary}

\begin{proof}
The map
\[
 U_M:=L_M\mathbf D_M^{\dagger/2}|_{S_M}
\]
is an isometry because \(U_M^*U_M=I_{S_M}\).  Also
\(\|C_M\|^2=d_{\rm miss}\) by the exact fiberwise Cauchy--Schwarz
calculation of TPC--60.  Hence
\[
 \kappa_M
 =\|C_MU_M\|^2
 \le\|C_M\|^2=d_{\rm miss}.
\]
Sharp examples are given in
\cref{ex:tpc61-coherent-incidence,ex:tpc61-fourier-incidence}.
\end{proof}

\subsection{Prime blocks and the unequal-prime alias operator}

For \(p>y\), let \(\cW_{M,p}\) be the missing atoms with terminal prime
\(p\), and let \(M_p\) be their contribution to the native map.  Thus
\begin{equation}
 M_X=\sum_{p>y}M_p.
 \label{eq:tpc61-prime-map-sum}
\end{equation}
Let
\begin{equation}
 \mathbf D_p
 :=\operatorname{diag}\left(
  \sum_{\substack{w\in\cW_{M,p}\\m(w)=m}}|\beta_w|^2
 \right)_{m\in\cM_X}.
 \label{eq:tpc61-prime-diagonal}
\end{equation}
The critical-cutoff fixed-\((i,p)\) singleton from TPC--59 and TPC--60
states that a fixed native output \(i\) and prime \(p\) support at most one
active high atom \citep{WangTPC59,WangTPC60}.

\begin{theorem}[Prime-block isometry and exact alias Gram]
\label{thm:tpc61-prime-block-identity}
Assume \(S_M\ne\{0\}\).  On the literal fixed-\(h_0\) carrier,
\begin{align}
 M_p^*M_p&=\mathbf D_p,
 \label{eq:tpc61-prime-block-isometry}\\
 \sum_{p>y}\mathbf D_p&=\mathbf D_M.
 \label{eq:tpc61-prime-diagonal-sum}
\end{align}
Define
\begin{equation}
 \widehat M_p
 :=M_p\mathbf D_M^{\dagger/2}|_{S_M},
 \qquad
 \Xi_{\ne p}
 :=\sum_{\substack{p,q>y\\p\ne q}}
     \widehat M_p^*\widehat M_q.
 \label{eq:tpc61-unequal-prime-operator}
\end{equation}
Then \(\Xi_{\ne p}\) is Hermitian and
\begin{align}
 \sum_p\widehat M_p^*\widehat M_p&=I_{S_M},
 \label{eq:tpc61-normalized-prime-resolution}\\
 \mathbf D_M^{\dagger/2}M_X^*M_X
 \mathbf D_M^{\dagger/2}\big|_{S_M}
 &=I_{S_M}+\Xi_{\ne p},
 \label{eq:tpc61-native-gram-prime-split}\\
 \kappa_M&=\boxed{1+\lambda_{\max}(\Xi_{\ne p})}.
 \label{eq:tpc61-kappa-unequal-prime}
\end{align}
Every nonzero matrix entry of \(\Xi_{\ne p}\) is an unequal-prime alias
at a common native output.
\end{theorem}

\begin{proof}
For a fixed \(p\), the singleton property makes the summands at distinct
atoms occupy distinct native coordinates.  Therefore
\[
 \|M_p\zeta\|^2
 =\sum_{w\in\cW_{M,p}}
   |\beta_w|^2|\zeta_{m(w)}|^2
 =\langle\mathbf D_p\zeta,\zeta\rangle,
\]
which proves \eqref{eq:tpc61-prime-block-isometry}.  The prime carrier
partition proves \eqref{eq:tpc61-prime-diagonal-sum}.  Normalize this
identity by \(\mathbf D_M^{\dagger/2}\) to obtain
\eqref{eq:tpc61-normalized-prime-resolution}.  Expanding the square of
\(\sum_p\widehat M_p\) gives
\eqref{eq:tpc61-native-gram-prime-split}; the ordered terms
\((p,q)\) and \((q,p)\) show that \(\Xi_{\ne p}\) is Hermitian.
Equation \eqref{eq:tpc61-kappa-unequal-prime} follows from
\cref{thm:tpc61-optimal-kappa}.
\end{proof}

The operator \(\Xi_{\ne p}\) is generally indefinite; only
\(I+\Xi_{\ne p}\) is positive.  This causes no loss in
\eqref{eq:tpc61-kappa-unequal-prime}: adding the identity shifts every
eigenvalue by exactly one, so
\(\lambda_{\max}(I+\Xi_{\ne p})
=1+\lambda_{\max}(\Xi_{\ne p})\) without any positivity assumption on
\(\Xi_{\ne p}\).

The diagonal of \(\Xi_{\ne p}\) is zero.  Indeed, \(i\) fixes the orbit
time \(j\), while \(i\) and a source row \(m\) fix
\(n=mj+h_0\), its unique prime \(p=\Pplus(n)>y\), the cofactor \(r=n/p\),
and the residual label \(s=d_mr\).  Thus two different terminal primes
from one source row cannot meet at the same native output.  In particular,
if \(r_M:=\dim S_M>0\), the normalized native Gram has trace \(r_M\).

\subsection{Same-residual and unequal-residual spectra}

Let \(T_{M,p}\) denote the contribution of \(\cW_{M,p}\) to the residual
map \(T_M\), and set
\begin{equation}
 \widehat T_p
 :=T_{M,p}\mathbf D_M^{\dagger/2}|_{S_M}.
 \label{eq:tpc61-normalized-residual-prime-map}
\end{equation}
Because the residual Hilbert space keeps \((i,s)\) orthogonal, the
cross-prime products in
\[
 \sum_{p\ne q}\widehat T_p^*\widehat T_q
\]
retain only collisions with equal residual label.

\begin{definition}[The two alias operators]
\label{def:tpc61-two-alias-operators}
Define
\begin{align}
 \Xi_{=s}
 &:=\sum_{\substack{p,q>y\\p\ne q}}
       \widehat T_p^*\widehat T_q,
 \label{eq:tpc61-same-residual-alias}\\
 \Xi_{\ne s}
 &:=\Xi_{\ne p}-\Xi_{=s}.
 \label{eq:tpc61-unequal-residual-alias}
\end{align}
\end{definition}

\begin{theorem}[Residual/native positive-spectrum split]
\label{thm:tpc61-residual-native-split}
The two operators in
\cref{def:tpc61-two-alias-operators} are Hermitian, have zero diagonal,
and satisfy
\begin{align}
 \mathbf D_M^{\dagger/2}T_M^*T_M
 \mathbf D_M^{\dagger/2}\big|_{S_M}
 &=I+\Xi_{=s},
 \label{eq:tpc61-residual-normalized-gram}\\
 \mathbf D_M^{\dagger/2}M_X^*M_X
 \mathbf D_M^{\dagger/2}\big|_{S_M}
 &=I+\Xi_{=s}+\Xi_{\ne s}.
 \label{eq:tpc61-native-normalized-gram-split}
\end{align}
Consequently, if
\begin{equation}
 \kappa_{\rm res}
 :=\sup_{\cD_M(\zeta)>0}
 \frac{\|T_M\zeta\|^2}{\cD_M(\zeta)},
 \label{eq:tpc61-residual-effective-degree}
\end{equation}
then
\begin{align}
 \kappa_{\rm res}
 &=1+\lambda_{\max}(\Xi_{=s}),
 \label{eq:tpc61-residual-positive-spectrum}\\
 \boxed{
 \kappa_M
 \le1+\lambda_{\max}(\Xi_{=s})
       +\lambda_{\max}(\Xi_{\ne s}).}
 \label{eq:tpc61-two-positive-spectrum-bound}
\end{align}
\end{theorem}

\begin{proof}
The fixed-\((i,s,p)\) singleton gives
\(T_{M,p}^*T_{M,p}=\mathbf D_p\).  Expanding
\(T_M^*T_M\), normalizing, and using
\eqref{eq:tpc61-prime-diagonal-sum} proves
\eqref{eq:tpc61-residual-normalized-gram}.  Subtract it from
\eqref{eq:tpc61-native-gram-prime-split} to obtain
\eqref{eq:tpc61-native-normalized-gram-split}.  Rayleigh--Ritz gives
\eqref{eq:tpc61-residual-positive-spectrum}, while the maximum-eigenvalue
form of Weyl's inequality gives
\eqref{eq:tpc61-two-positive-spectrum-bound}.
\end{proof}

The arithmetic roles of the two terms differ.  The support of
\(\Xi_{=s}\) is governed by the determinant ladder
\[
 rp-dj\ell=h_0
\]
inside one residual label.  The operator \(\Xi_{\ne s}\) is created only
after \(s\) is erased; its terms share the native output \(i\) but have
different residual labels.  Thus the residual ladder does not by itself
control \(\Xi_{\ne s}\).  Also, the present upper-norm problem needs only
the positive spectra of these signed operators.  Bounding their full
absolute norms is sufficient but can discard favorable cancellation.

\subsection{The quotient on the actual residual image}

Put
\begin{equation}
 \mathbf R_M:=T_M^*T_M,
 \qquad
 \mathbf N_M:=T_M^*\mathsf E^*\mathsf E T_M=M_X^*M_X,
 \qquad
 S_R:=\Ran\mathbf R_M=(\Ker T_M)^\perp,
 \label{eq:tpc61-residual-native-gram-pair}
\end{equation}
and let
\begin{equation}
 \mathcal E_M:=\Ran T_M
 \label{eq:tpc61-actual-residual-image}
\end{equation}
be the actual residual coefficient image.

\begin{theorem}[Exact restricted erasure norm]
\label{thm:tpc61-restricted-erasure-norm}
If \(\mathcal E_M\ne\{0\}\), then
\begin{align}
 \beta_M
 &:=\sup_{0\ne z\in\mathcal E_M}
   \frac{\|\mathsf Ez\|^2}{\|z\|^2}
 \label{eq:tpc61-beta-definition}\\
 &=\lambda_{\max}\!\left(
   \mathbf R_M^{\dagger/2}\mathbf N_M
   \mathbf R_M^{\dagger/2}\big|_{S_R}
   \right).
 \label{eq:tpc61-beta-quotient-formula}
\end{align}
Moreover,
\begin{equation}
 \boxed{\kappa_M\le\kappa_{\rm res}\beta_M.}
 \label{eq:tpc61-two-stage-upper}
\end{equation}
\end{theorem}

\begin{proof}
The polar-normalized map
\[
 V:=T_M\mathbf R_M^{\dagger/2}|_{S_R}:S_R\longrightarrow\mathcal E_M
\]
is unitary.  Conjugating
\(\mathsf E^*\mathsf E|_{\mathcal E_M}\) by \(V\) gives the operator in
\eqref{eq:tpc61-beta-quotient-formula}.  Finally,
\[
 \|M_X\zeta\|^2
 \le\beta_M\|T_M\zeta\|^2
 \le\beta_M\kappa_{\rm res}\cD_M(\zeta),
\]
and optimality in \cref{thm:tpc61-optimal-kappa} proves
\eqref{eq:tpc61-two-stage-upper}.
\end{proof}

For comparison with the coefficient-blind ladder bounds, put
\begin{align}
 d_{\rm res}
 &:=\max_{i,s}
   |\{w\in\cW_M:i(w)=i,\ s(w)=s\}|,
 \label{eq:tpc61-ambient-residual-degree}\\
 k_{\rm erase}
 &:=\max_i
   |\{s:(i,s)\text{ is active in the missing residual carrier}\}|.
 \label{eq:tpc61-ambient-erasure-degree}
\end{align}
The two ambient grouping maps have squared norms \(d_{\rm res}\) and
\(k_{\rm erase}\), respectively.  Compression to the physical images
therefore gives
\begin{equation}
 \boxed{
 \kappa_{\rm res}\le d_{\rm res},
 \qquad
 \beta_M\le k_{\rm erase}.}
 \label{eq:tpc61-two-ambient-comparisons}
\end{equation}
Neither reverse comparison is valid: the common row coefficients and the
literal phases may make the restricted norms much smaller.

The ambient erasure norm is the largest number of active residual labels
at one \(i\), but \(\beta_M\) is its compression to
\(\mathcal E_M\).  These need not be comparable from below.  Likewise,
the product in \eqref{eq:tpc61-two-stage-upper} is a useful certificate,
not an identity in general; the one-stage constant \(\kappa_M\) remains
optimal.

\begin{remark}[Block notation and the residual M\"obius gauge]
\label{rem:tpc61-block-erasure-notation}
On a dyadic cofactor block \(R\), the quantity \(\beta_{M,R}\) in
\eqref{eq:tpc61-restricted-erasure-exponent} is precisely the specialization
of \(\beta_M\) to the corresponding atomic and residual maps.  If the
residual M\"obius sign is instead stored in a final signed erasure
\(C_\mu\), rather than in the residual input as in
\eqref{eq:tpc61-sign-compatible-atom}, the two descriptions are conjugate
by the diagonal unitary
\[
 (U_\mu z)_{i,s}=\mu(s)z_{i,s}.
\]
Thus the two gauges have the same restricted norm.  They describe one gate
and must not both be charged in
the endpoint ledger.
\end{remark}

\subsection{Uniform and distinguished quantifiers}

For one coefficient vector with \(\cD_M(\zeta)>0\), define
\begin{equation}
 \kappa_M(\zeta)
 :=\frac{\|M_X\zeta\|^2}{\cD_M(\zeta)}.
 \label{eq:tpc61-pointwise-kappa}
\end{equation}

\begin{proposition}[Pointwise Rayleigh identity]
\label{prop:tpc61-pointwise-rayleigh}
Let \(x=\mathbf D_M^{1/2}\zeta\).  Then
\begin{equation}
 \boxed{
 \kappa_M(\zeta)
 =\frac{\langle(I+\Xi_{\ne p})x,x\rangle}{\|x\|^2},
 \qquad
 0\le\kappa_M(\zeta)\le\kappa_M.}
 \label{eq:tpc61-pointwise-rayleigh}
\end{equation}
The pointwise value may be zero even when the uniform constant is large.
\end{proposition}

\begin{proof}
Insert \(x=\mathbf D_M^{1/2}\zeta\) into
\eqref{eq:tpc61-native-gram-prime-split}.  Positivity gives the lower
bound and \cref{thm:tpc61-optimal-kappa} gives the upper bound.
\end{proof}

For the represented map \(A_X\) and every vector with
\(A_X\zeta\ne0\), the reverse triangle inequality now gives the optimal
mass-to-native certificate available from \(\cD_M\):
\begin{align}
 \frac{\|(A_X+M_X)\zeta\|^2}{\|A_X\zeta\|^2}
 &\ge
 \left(
  1-\sqrt{
   \frac{\kappa_M\cD_M(\zeta)}
        {\|A_X\zeta\|^2}}
 \right)_+^2
 \label{eq:tpc61-uniform-reassembly-certificate}\\
 &\ge
 \left(
  1-\sqrt{
   \frac{d_{\rm miss}\cD_M(\zeta)}
        {\|A_X\zeta\|^2}}
 \right)_+^2.
 \notag
\end{align}
For one distinguished physical vector, \(\kappa_M\) in the first line can
be replaced by the exact number \(\kappa_M(\zeta)\).  This replacement is
pointwise only: a small Rayleigh quotient for one vector does not imply the
Loewner inequality \eqref{eq:tpc61-kappa-loewner}.
